% Part: first-order-logic
% Chapter: completeness
% Section: lindenbaums-lemma

\documentclass[../../../include/open-logic-section]{subfiles}

\begin{document}

\iftag{FOL}
      {\olfileid{fol}{com}{lin}}
      {\olfileid{pl}{com}{lin}}

\olsection{లిండెన్‌బామ్ ఉపసిద్ధాంతం}

\begin{explain}
ఏ అవైరుధ్య !!{sentence}s సమితైనా కేవలం అవైరుధ్యమైనదే కాక
!!{complete} కూడా అయిన ఏదో ఒక వాక్యాల సమితిలో ఇమిడి ఉంటుందని చూపే
ఉపసిద్ధాంతాన్ని ఇప్పుడు నిరూపిస్తాం. సమితి అవైరుధ్యంగానే ఉండటానికి ప్రతి
అడుగులో హామీ ఇస్తూ, ఒకసారి ఒక !!{sentence}ను చేర్చడం ద్వారా నిరూపణ
పనిచేస్తుంది. ప్రతి~$!A$కు ఏదో ఒక దశలో $!A$ లేదా $\lnot !A$
చేరేలా దీన్ని చేస్తాం. అప్పుడు ఆ నిర్మాణంలోని అన్ని దశల సమ్మేళనంలో
$!A$ లేదా దాని నిషేధం~$\lnot !A$ ఉంటుంది; అందువల్ల అది సంపూర్ణం.
ప్రతి దశలోనూ వైరుధ్యాన్ని ప్రవేశపెట్టకుండా చూసుకుంటాం కనుక అది
అవైరుధ్యమైనది కూడా.
\end{explain}

\begin{lem}[లిండెన్‌బామ్ ఉపసిద్ధాంతం]
\ollabel{lem:lindenbaum} భాష~$\Lang{L}$లోని ప్రతి అవైరుధ్య
సమితి~$\Gamma$ను !!a{complete}, అవైరుధ్య సమితి~$\Gamma^*$గా
విస్తరించవచ్చు.
\end{lem}

\begin{proof}
$\Gamma$ అవైరుధ్యమని అనుకుందాం. $\Lang L$లోని అన్ని !!{sentence}sకు
$!A_0$, $!A_1$, \dots{} అనే ఒక లెక్కింపు ఉండనివ్వండి. $\Gamma_0 =
\Gamma$గా, ఇంకా
\[
\Gamma_{n+1} =
\begin{cases}
\Gamma_n \cup \{ !A_n \} & \textrm{if $\Gamma_n \cup \{!A_n\}$ is
  consistent;} \\
\Gamma_n \cup \{ \lnot !A_n \} & \textrm{otherwise.}
\end{cases}
\]
గా నిర్వచించండి. $\Gamma^* = \bigcup_{n \geq 0} \Gamma_n$గా ఉంచండి.

ప్రతి~$\Gamma_n$ అవైరుధ్యమైనది: నిర్వచనం ప్రకారం $\Gamma_0$
అవైరుధ్యమైనది. $\Gamma_{n+1} = \Gamma_n \cup \{!A_n\}$ అయితే,
చివరిది అవైరుధ్యమైనందువల్లే అలా ఉంటుంది. అది అవైరుధ్యం కాకపోతే,
$\Gamma_{n+1} = \Gamma_n \cup \{\lnot !A_n\}$. $\Gamma_n \cup
\{\lnot !A_n\}$ అవైరుధ్యమని నిర్ధారించాలి. అది కూడా అవైరుధ్యం కాదని
అనుకుందాం. అప్పుడు $\Gamma_n \cup \{!A_n\}$, $\Gamma_n \cup
\{\lnot !A_n\}$
\emph{రెండూ} వైరుధ్యభరితమైనవి. కింది ఫలితం ప్రకారం ఇది~$\Gamma_n$
వైరుధ్యభరితమని అర్థం:
\iftag{FOL}{%
  \tagrefs{prfAX/{fol:axd:prv:prop:provability-exhaustive},
    prfSC/{fol:seq:prv:prop:provability-exhaustive},
    prfND/{fol:ntd:prv:prop:provability-exhaustive},
    prfTab/{fol:tab:prv:prop:provability-exhaustive}}}{%
  \tagrefs{prfAX/{pl:axd:prv:prop:provability-exhaustive},
    prfSC/{pl:seq:prv:prop:provability-exhaustive},
    prfND/{pl:ntd:prv:prop:provability-exhaustive},
    prfTab/{pl:tab:prv:prop:provability-exhaustive}}};
ఇది ఆగమన పరికల్పనకు విరుద్ధం.

ప్రతి~$n$కు, ప్రతి $i < n$కు $\Gamma_i \subseteq \Gamma_n$. ఇది
~$n$పై సరళమైన ఆగమనం ద్వారా వస్తుంది. $n=0$ అయినప్పుడు $i < 0$
అయ్యేదేమీ లేదు; కాబట్టి వాదన స్వయంచాలకంగా నిజం. ఆగమన అడుగులో అది
~$n$కు నిజమని అనుకుందాం. $i < n+1$ అయితే $\Gamma_i \subseteq
\Gamma_{n+1}$ అని చూపుతాం. నిర్మాణం ప్రకారం $\Gamma_{n+1} = \Gamma_n
\cup \{!A_n\}$ లేదా $= \Gamma_n \cup \{\lnot !A_n\}$. కనుక
$\Gamma_n \subseteq \Gamma_{n+1}$. $i < n+1$ అయితే, ఆగమన పరికల్పన
ప్రకారం ($i < n$ అయితే) $\Gamma_i \subseteq \Gamma_n$; లేదా ($i = n$
అయితే) $\Gamma_n \subseteq \Gamma_n$ అనే తక్షణ వాస్తవం వర్తిస్తుంది.
$\subseteq$ యొక్క సంక్రామకత్వం ద్వారా $\Gamma_i \subseteq
\Gamma_{n+1}$ వస్తుంది.

దీని నుంచి $\Gamma^*$ అవైరుధ్యమని అనుగమిస్తుంది. కారణం ఇది: పరిమిత
$\Gamma' \subseteq \Gamma^*$ను తీసుకోండి. ప్రతి $!B \in \Gamma'$
ఏదో ఒక~$i$కు~$\Gamma_i$లో కూడా ఉంటుంది. వాటిలో అతిపెద్దదాన్ని $n$
అందాం. $i \le n$ అయితే $\Gamma_i \subseteq \Gamma_n$ కనుక, ప్రతి
$!B \in \Gamma'$కు $!B \in \Gamma_n$ కూడా; అంటే $\Gamma' \subseteq
\Gamma_n$, పైగా $\Gamma_n$ అవైరుధ్యమైనది.
\sourcecorrection{OLTECOM-009}{ఆంగ్ల మూలం ఇక్కడ ఎడమ సభ్యమే లేని
$\in \Gamma_n$ను రాసింది. వాక్యం పరిధిలో ఉన్న $!B$ను పునరుద్ధరించి
పూర్తి సభ్యత్వ ప్రకటనగా రాశాం.} అందువల్ల ప్రతి పరిమిత
ఉపసమితి $\Gamma' \subseteq \Gamma^*$ అవైరుధ్యమైనది. కింది ఫలితం
ప్రకారం $\Gamma^*$ అవైరుధ్యమైనది:
\iftag{FOL}{%
  \tagrefs{prfAX/{fol:axd:ptn:prop:proves-compact},
    prfSC/{fol:seq:ptn:prop:proves-compact},
    prfND/{fol:ntd:ptn:prop:proves-compact},
    prfTab/{fol:tab:ptn:prop:proves-compact}}}{%
  \tagrefs{prfAX/{pl:axd:ptn:prop:proves-compact},
    prfSC/{pl:seq:ptn:prop:proves-compact},
    prfND/{pl:ntd:ptn:prop:proves-compact},
    prfTab/{pl:tab:ptn:prop:proves-compact}}}.

$\Frm[L]$లోని ప్రతి !!{sentence} $\Gamma^*$ను నిర్వచించడానికి
ఉపయోగించిన జాబితాలో కనిపిస్తుంది. $!A_n \notin \Gamma^*$ అయితే,
$\Gamma_n \cup \{!A_n\}$ వైరుధ్యభరితమై ఉండటమే దానికి కారణం. కానీ
అప్పుడు $\lnot !A_n \in \Gamma^*$; అందువల్ల $\Gamma^*$
!!{complete}.
\end{proof}

\end{document}
